\documentclass[11pt,a4paper]{article}
\usepackage[top=2.3cm,bottom=2.3cm,left=2.2cm,right=2.2cm]{geometry}

\usepackage{amsmath,amssymb}
\usepackage{fontspec}
\setmainfont{Liberation Serif}
\setsansfont{Liberation Sans}
\setmonofont{DejaVu Sans Mono}[Scale=0.84]

\usepackage{xcolor}
\definecolor{brandblue}{RGB}{20, 55, 105}
\definecolor{accentblue}{RGB}{35, 95, 165}
\definecolor{codebg}{RGB}{248, 249, 251}
\definecolor{codeframe}{RGB}{215, 222, 230}
\definecolor{javakeyword}{RGB}{0, 45, 170}
\definecolor{javastring}{RGB}{5, 120, 20}
\definecolor{javacomment}{RGB}{120, 120, 120}
\definecolor{javanumber}{RGB}{20, 75, 220}

\definecolor{noteborder}{RGB}{30, 110, 65}
\definecolor{notebg}{RGB}{242, 249, 245}
\definecolor{warnborder}{RGB}{185, 75, 10}
\definecolor{warnbg}{RGB}{254, 248, 240}
\definecolor{tipborder}{RGB}{30, 95, 160}
\definecolor{tipbg}{RGB}{242, 247, 254}

\usepackage{fancyhdr}
\usepackage{lastpage}
\pagestyle{fancy}
\fancyhf{}
\lhead{\parbox[t]{0.58\textwidth}{\raggedright\small\sffamily\color{brandblue}\textbf{T18 --- Interfacce Funzionali, Espressioni Lambda e Streams API}}}
\rhead{\small\sffamily\color{gray}Programmazione II -- UniVR (2025/2026)}
\cfoot{\small\sffamily Pagina \thepage\ di \pageref{LastPage}}
\renewcommand{\headrulewidth}{0.5pt}
\renewcommand{\headrule}{\hbox to\headwidth{\color{brandblue!70!black}\leaders\hrule height \headrulewidth\hfill}}

\usepackage{titlesec}
\titleformat{\section}{\Large\bfseries\sffamily\color{brandblue}}{\thesection}{0.8em}{}[\titlerule]
\titleformat{\subsection}{\large\bfseries\sffamily\color{accentblue}}{\thesubsection}{0.8em}{}
\titleformat{\subsubsection}{\normalsize\bfseries\sffamily\color{brandblue!85!black}}{\thesubsubsection}{0.8em}{}

\usepackage{needspace}
\usepackage{fancyvrb}
\DefineVerbatimEnvironment{asciibox}{Verbatim}{
  fontsize=\fontsize{7.5pt}{9.2pt}\selectfont,
  frame=single,
  rulecolor=\color{codeframe},
  fillcolor=\color{codebg},
  framesep=2.5mm
}
\usepackage{listings}
\lstset{
  backgroundcolor=\color{codebg},
  frame=single,
  rulecolor=\color{codeframe},
  basicstyle=\ttfamily\footnotesize,
  keywordstyle=\color{javakeyword}\bfseries,
  stringstyle=\color{javastring},
  commentstyle=\color{javacomment}\itshape,
  numberstyle=\tiny\color{gray},
  numbers=left,
  numbersep=7pt,
  stepnumber=1,
  breaklines=true,
  breakatwhitespace=true,
  showstringspaces=false,
  tabsize=4,
  xleftmargin=12pt,
  framexleftmargin=12pt
}

\newcommand{\passthrough}[1]{#1}
\providecommand{\tightlist}{\setlength{\itemsep}{0pt}\setlength{\parskip}{0pt}}

\usepackage{tcolorbox}
\tcbuselibrary{skins,breakable}
\newtcolorbox{studynote}[1][Nota]{
  enhanced,
  breakable,
  colback=notebg,
  colframe=noteborder,
  fonttitle=\bfseries\sffamily\small,
  coltitle=white,
  title={#1},
  arc=2mm,
  boxrule=0.8pt,
  left=9pt, right=9pt, top=7pt, bottom=7pt
}

\newtcolorbox{studywarning}[1][Attenzione]{
  enhanced,
  breakable,
  colback=warnbg,
  colframe=warnborder,
  fonttitle=\bfseries\sffamily\small,
  coltitle=white,
  title={#1},
  arc=2mm,
  boxrule=0.8pt,
  left=9pt, right=9pt, top=7pt, bottom=7pt
}

\newtcolorbox{studytip}[1][Suggerimento]{
  enhanced,
  breakable,
  colback=tipbg,
  colframe=tipborder,
  fonttitle=\bfseries\sffamily\small,
  coltitle=white,
  title={#1},
  arc=2mm,
  boxrule=0.8pt,
  left=9pt, right=9pt, top=7pt, bottom=7pt
}

\usepackage{booktabs}
\usepackage{longtable}
\usepackage{array}
\usepackage{calc}

\usepackage{hyperref}
\hypersetup{
  colorlinks=true,
  linkcolor=accentblue,
  urlcolor=accentblue,
  citecolor=brandblue,
  pdfborder={0 0 0}
}

\title{\Huge\bfseries\sffamily\color{brandblue} T18 --- Interfacce Funzionali, Espressioni Lambda e Streams API \\[0.3em] \large\sffamily Predicate, Function, Consumer, Supplier, Method Reference (::) e Pipeline Stream}
\author{\textbf{Docente:} Prof. Michele Pasqua \quad---\quad \textbf{Appunti Revisione:} Wally}
\date{A.A. 2025/2026 -- Universit\`a degli Studi di Verona}

\begin{document}
\maketitle
\thispagestyle{fancy}
\vspace{-1.2em}
\noindent\textcolor{brandblue}{\rule{\textwidth}{0.8pt}}
\vspace{1.2em}

\begin{center}\rule{0.5\linewidth}{0.5pt}\end{center}

\subsection{1. Analisi Teorica Approfondita (Slide
T18)}\label{analisi-teorica-approfondita-slide-t18}

La lezione 18 introduce il paradigma di \textbf{programmazione
funzionale e dichiarativa} in Java (introdotto in Java 8), imperniato su
due pilastri: 1. \textbf{Interfacce Funzionali} e \textbf{Method
Reference} (\passthrough{\lstinline!::!}). 2. \textbf{Stream API}
(\passthrough{\lstinline!java.util.stream!}), per l'elaborazione fluente
e parallela di flussi di dati.

\begin{center}\rule{0.5\linewidth}{0.5pt}\end{center}

\subsubsection{1.1 Cos'è una Functional Interface? (Slide
3)}\label{cosuxe8-una-functional-interface-slide-3}

Un'interfaccia si dice \textbf{funzionale} se dichiara
\textbf{esattamente un unico metodo astratto} (pattern noto come
\textbf{SAM} --- \emph{Single Abstract Method}). * \textbf{Semantica}:
Puramente funzionale; l'esito della chiamata dipende esclusivamente
dagli argomenti ricevuti (nessun effetto collaterale occulto). *
\textbf{L'annotazione \passthrough{\lstinline!@FunctionalInterface!}}: *
Non è obbligatoria per far funzionare una lambda, ma è una best practice
fondamentale: istruisce il compilatore Java a verificare che
l'interfaccia contenga un solo metodo astratto. Se qualcuno tenta di
aggiungere un secondo metodo astratto, il compilatore rigetta il codice
con errore. * \textbf{Metodi esclusi dal conteggio SAM}: * Metodi con
implementazione predefinita (\passthrough{\lstinline!default!}). *
Metodi statici (\passthrough{\lstinline!static!}). * Metodi pubblici
astratti che fanno override di metodi di
\passthrough{\lstinline!java.lang.Object!} (es.
\passthrough{\lstinline!boolean equals(Object obj)!}).

\begin{center}\rule{0.5\linewidth}{0.5pt}\end{center}

\subsubsection{\texorpdfstring{1.2 Le 4 Interfacce Funzionali
Fondamentali (\texttt{java.util.function}) (Slide
4-12)}{1.2 Le 4 Interfacce Funzionali Fondamentali (java.util.function) (Slide 4-12)}}\label{le-4-interfacce-funzionali-fondamentali-java.util.function-slide-4-12}

Il package \passthrough{\lstinline!java.util.function!} standardizza i
contratti funzionali più frequenti:

\begin{longtable}[]{@{}
  >{\raggedright\arraybackslash}p{(\linewidth - 6\tabcolsep) * \real{0.2500}}
  >{\raggedright\arraybackslash}p{(\linewidth - 6\tabcolsep) * \real{0.2500}}
  >{\raggedright\arraybackslash}p{(\linewidth - 6\tabcolsep) * \real{0.2500}}
  >{\raggedright\arraybackslash}p{(\linewidth - 6\tabcolsep) * \real{0.2500}}@{}}
\toprule\noalign{}
\begin{minipage}[b]{\linewidth}\raggedright
Interfaccia
\end{minipage} & \begin{minipage}[b]{\linewidth}\raggedright
Firma Metodo SAM
\end{minipage} & \begin{minipage}[b]{\linewidth}\raggedright
Semantica Funzionale
\end{minipage} & \begin{minipage}[b]{\linewidth}\raggedright
Esempio d'Uso Tipico
\end{minipage} \\
\midrule\noalign{}
\endhead
\bottomrule\noalign{}
\endlastfoot
\passthrough{\lstinline!Predicate<T>!} &
\passthrough{\lstinline!boolean test(T t)!} & Filtro / Condizione
booleana & \passthrough{\lstinline!n -> n \% 2 == 0!} \\
\passthrough{\lstinline!Function<T, R>!} &
\passthrough{\lstinline!R apply(T t)!} & Trasformazione \(T \to R\) &
\passthrough{\lstinline!str -> str.length()!} \\
\passthrough{\lstinline!Consumer<T>!} &
\passthrough{\lstinline!void accept(T t)!} & Azione con effetto
collaterale &
\passthrough{\lstinline!str -> System.out.println(str)!} \\
\passthrough{\lstinline!Supplier<T>!} &
\passthrough{\lstinline!T get()!} & Generatore / Factory &
\passthrough{\lstinline!() -> Math.random()!} \\
\end{longtable}

\paragraph{Varianti e Specializzazioni (Slide
6-12)}\label{varianti-e-specializzazioni-slide-6-12}

\begin{enumerate}
\def\labelenumi{\arabic{enumi}.}
\tightlist
\item
  \textbf{Versioni Primitive}: Per evitare l'overhead di
  boxing/unboxing, Java fornisce interfacce specializzate per
  \passthrough{\lstinline!int!}, \passthrough{\lstinline!long!},
  \passthrough{\lstinline!double!}:

  \begin{itemize}
  \tightlist
  \item
    \passthrough{\lstinline!IntPredicate!}
    (\passthrough{\lstinline!boolean test(int v)!}),
    \passthrough{\lstinline!LongPredicate!},
    \passthrough{\lstinline!DoublePredicate!}.
  \item
    \passthrough{\lstinline!IntFunction<R>!},
    \passthrough{\lstinline!ToIntFunction<T>!},
    \passthrough{\lstinline!IntConsumer!},
    \passthrough{\lstinline!IntSupplier!}.
  \end{itemize}
\item
  \textbf{Varianti a Due Argomenti}:

  \begin{itemize}
  \tightlist
  \item
    \passthrough{\lstinline!BiPredicate<T, U>!}:
    \passthrough{\lstinline!boolean test(T t, U u)!} (es.
    \passthrough{\lstinline!(a, b) -> a > b!}).
  \item
    \passthrough{\lstinline!BiFunction<T, U, R>!}:
    \passthrough{\lstinline!R apply(T t, U u)!} (es.
    \passthrough{\lstinline!(a, b) -> a + b!}).
  \item
    \passthrough{\lstinline!BiConsumer<T, U>!}:
    \passthrough{\lstinline!void accept(T t, U u)!}.
  \end{itemize}
\item
  \textbf{Operatori (Input e Output dello stesso tipo)}:

  \begin{itemize}
  \tightlist
  \item
    \passthrough{\lstinline!UnaryOperator<T> extends Function<T, T>!}:
    riceve un \(T\) e restituisce un \(T\).
  \item
    \passthrough{\lstinline!BinaryOperator<T> extends BiFunction<T, T, T>!}:
    riceve due \(T\) e restituisce un \(T\).
  \end{itemize}
\end{enumerate}

\begin{center}\rule{0.5\linewidth}{0.5pt}\end{center}

\subsubsection{\texorpdfstring{1.3 I Quattro Tipi di Method Reference
(\texttt{::}) (Slide
13-15)}{1.3 I Quattro Tipi di Method Reference (::) (Slide 13-15)}}\label{i-quattro-tipi-di-method-reference-slide-13-15}

Il \emph{Method Reference} è uno zucchero sintattico introdotto in Java
8 per rendere le espressioni lambda ancora più concise quando queste si
limitano a inoltrare direttamente i propri parametri a un metodo
esistente:

\begin{longtable}[]{@{}
  >{\raggedright\arraybackslash}p{(\linewidth - 4\tabcolsep) * \real{0.3333}}
  >{\raggedright\arraybackslash}p{(\linewidth - 4\tabcolsep) * \real{0.3333}}
  >{\raggedright\arraybackslash}p{(\linewidth - 4\tabcolsep) * \real{0.3333}}@{}}
\toprule\noalign{}
\begin{minipage}[b]{\linewidth}\raggedright
Categoria
\end{minipage} & \begin{minipage}[b]{\linewidth}\raggedright
Sintassi Method Reference
\end{minipage} & \begin{minipage}[b]{\linewidth}\raggedright
Equivalente Lambda Esplicita
\end{minipage} \\
\midrule\noalign{}
\endhead
\bottomrule\noalign{}
\endlastfoot
\textbf{1. Metodo Statico} &
\passthrough{\lstinline!Math::random!}\passthrough{\lstinline!Integer::parseInt!}
&
\passthrough{\lstinline!() -> Math.random()!}\passthrough{\lstinline!str -> Integer.parseInt(str)!} \\
\textbf{2. Metodo d'Istanza su Oggetto Esistente} (\emph{Bound}) &
\passthrough{\lstinline!System.out::println!}\passthrough{\lstinline!hexDigits::charAt!}
&
\passthrough{\lstinline!x -> System.out.println(x)!}\passthrough{\lstinline!i -> hexDigits.charAt(i)!} \\
\textbf{3. Metodo d'Istanza su Tipo Arbitrario} (\emph{Unbound}) &
\passthrough{\lstinline!String::toUpperCase!}\passthrough{\lstinline!String::length!}
&
\passthrough{\lstinline!str -> str.toUpperCase()!}\passthrough{\lstinline!str -> str.length()!} \\
\textbf{4. Costruttore} &
\passthrough{\lstinline!Person::new!}\passthrough{\lstinline!Integer[]::new!}
&
\passthrough{\lstinline!str -> new Person(str)!}\passthrough{\lstinline!size -> new Integer[size]!} \\
\end{longtable}

\begin{center}\rule{0.5\linewidth}{0.5pt}\end{center}

\subsubsection{\texorpdfstring{1.4 La Stream API
(\texttt{java.util.stream}) (Slide
16-18)}{1.4 La Stream API (java.util.stream) (Slide 16-18)}}\label{la-stream-api-java.util.stream-slide-16-18}

Uno \passthrough{\lstinline!Stream<T>!} è una \textbf{sequenza di
elementi generata da una sorgente} che supporta operazioni aggregate di
calcolo.

\paragraph{Le Tre Proprietà Chiave di uno
Stream:}\label{le-tre-proprietuxe0-chiave-di-uno-stream}

\begin{enumerate}
\def\labelenumi{\arabic{enumi}.}
\tightlist
\item
  \textbf{Pipelining}: Le operazioni intermedie restituiscono un nuovo
  stream, permettendo la concatenazione fluente (\emph{method
  chaining}).
\item
  \textbf{Internal Iteration}: A differenza delle collezioni dove lo
  sviluppatore controlla l'iterazione tramite
  \passthrough{\lstinline!for!}/\passthrough{\lstinline!while!}
  (\emph{iterazione esterna}), lo stream gestisce il ciclo internamente.
  Ciò consente alla JVM di ottimizzare l'esecuzione e abilitare il
  parallelismo trasparente (\passthrough{\lstinline!parallelStream()!}).
\item
  \textbf{Lazy Evaluation (Valutazione Pigra)}: Le operazioni intermedie
  non vengono eseguite quando vengono dichiarate; restano ``in attesa''.
  L'intera pipeline viene elaborata in un unico passaggio solo quando
  viene invocata un'operazione terminale!
\end{enumerate}

\begin{center}\rule{0.5\linewidth}{0.5pt}\end{center}

\subsubsection{1.5 Ciclo di Vita di una Pipeline Stream (Slide
18-24)}\label{ciclo-di-vita-di-una-pipeline-stream-slide-18-24}

Una pipeline si compone tassativamente di 3 fasi:

\[\text{Sorgente (Source)} \longrightarrow \text{Zero o più Operazioni Intermedie} \longrightarrow \text{Una Operazione Terminale}\]

\needspace{7cm}
\begin{asciibox}
┌───────────────┐     ┌─────────────┐     ┌───────────┐     ┌───────────────┐
│ Stream.of(...) ├──► │ .filter(...) ├──► │ .map(...) ├──►  │ .forEach(...) │
└───────────────┘     └─────────────┘     └───────────┘     └───────────────┘
   [Sorgente]          [Intermedia]        [Intermedia]        [Terminale]
\end{asciibox}

\paragraph{1. Creazione della Sorgente (Slide 19,
21-22)}\label{creazione-della-sorgente-slide-19-21-22}

\begin{itemize}
\tightlist
\item
  Da array: \passthrough{\lstinline!Arrays.stream(arr)!} o
  \passthrough{\lstinline!Stream.of(v1, v2, v3)!}.
\item
  Da collezione: \passthrough{\lstinline!list.stream()!} o
  \passthrough{\lstinline!set.stream()!}.
\item
  \textbf{Generazione infinita}:

  \begin{itemize}
  \tightlist
  \item
    \passthrough{\lstinline!Stream.iterate(seed, UnaryOperator)!}:
    applica iterativamente l'operatore a partire dal seme. Es:
    \passthrough{\lstinline!Stream.iterate(0, i -> i + 1)!} genera
    \(0, 1, 2, 3, \dots\)
  \item
    \passthrough{\lstinline!Stream.generate(Supplier)!}: invoca
    continuamente il fornitore. Es:
    \passthrough{\lstinline!Stream.generate(Math::random)!}.
  \item
    Per evitare loop infiniti, si arresta la generazione con
    l'operazione intermedia \passthrough{\lstinline!.limit(n)!}.
  \end{itemize}
\end{itemize}

\paragraph{2. Operazioni Intermedie (Slide 23, 25-29,
34-36)}\label{operazioni-intermedie-slide-23-25-29-34-36}

Restituiscono un nuovo \passthrough{\lstinline!Stream<R>!} e sono pigre
(\emph{lazy}): * \passthrough{\lstinline!filter(Predicate<T>)!}:
trattiene solo gli elementi per cui il predicato è
\passthrough{\lstinline!true!}. *
\passthrough{\lstinline!map(Function<T, R>)!}: trasforma ciascun
elemento \(T \to R\) (mappatura \(1 \to 1\)). *
\passthrough{\lstinline!flatMap(Function<T, Stream<R>>)!}: appiattisce
strutture annidate (es. matrici \passthrough{\lstinline!Integer[][]!} o
liste di liste) in un unico stream monodimensionale. *
\passthrough{\lstinline!distinct()!}: elimina i duplicati (in base a
\passthrough{\lstinline!equals!} e \passthrough{\lstinline!hashCode!}).
* \passthrough{\lstinline!sorted()!} /
\passthrough{\lstinline!sorted(Comparator)!}: ordina gli elementi. *
\passthrough{\lstinline!limit(n)!}: tronca lo stream ai primi \(n\)
elementi. * \passthrough{\lstinline!skip(n)!}: scarta i primi \(n\)
elementi.

\paragraph{3. Operazioni Terminali (Slide 20, 23,
30-33)}\label{operazioni-terminali-slide-20-23-30-33}

Consumano lo stream e producono un risultato o un effetto collaterale: *
\textbf{Chiusura su Collezioni / Array (Slide 20)}: *
\passthrough{\lstinline!stream.toList()!}: raccoglie gli elementi in una
\passthrough{\lstinline!List!} immutabile (Java 16+). *
\passthrough{\lstinline!stream.collect(Collectors.toSet())!}: raccoglie
in un \passthrough{\lstinline!Set!}. *
\passthrough{\lstinline!stream.toArray(Integer[]::new)!}: esporta in
array tipizzato. * \textbf{Iterazione / Consumo}:
\passthrough{\lstinline!stream.forEach(Consumer<T>)!}. *
\textbf{Conteggio / Riduzione}:
\passthrough{\lstinline!stream.count()!},
\passthrough{\lstinline!stream.reduce(...)!}. * \textbf{Matching}:
\passthrough{\lstinline!anyMatch(Predicate)!},
\passthrough{\lstinline!allMatch(Predicate)!},
\passthrough{\lstinline!noneMatch(Predicate)!}. * \textbf{Ricerca \&
\passthrough{\lstinline!Optional<T>!} (Slide 31-32)}: *
\passthrough{\lstinline!findFirst()!},
\passthrough{\lstinline!findAny()!}: restituiscono un
\passthrough{\lstinline!Optional<T>!}, contenitore sicuro introdotto per
debellare le \passthrough{\lstinline!NullPointerException!}. * Metodi di
\passthrough{\lstinline!Optional!}:
\passthrough{\lstinline!isPresent()!}, \passthrough{\lstinline!get()!},
\passthrough{\lstinline!orElse(valoreDefault)!}.

\begin{studywarning}[Attenzione / Trappola d'Esame] \textbf{Regola Aurea degli Stream (Slide 24 --- Domanda
d'Esame):} \textbf{Uno Stream può essere consumato UNA SOLA VOLTA!} Se
si tenta di invocare una seconda operazione terminale sullo stesso
oggetto \passthrough{\lstinline!Stream!}, la JVM solleva a runtime:
\passthrough{\lstinline!java.lang.IllegalStateException: stream has already been operated upon or closed!}.
\end{studywarning}

\begin{center}\rule{0.5\linewidth}{0.5pt}\end{center}

\subsection{\texorpdfstring{2. Analisi Dettagliata del Codice
(\texttt{Lezione17/MavenDate})}{2. Analisi Dettagliata del Codice (Lezione17/MavenDate)}}\label{analisi-dettagliata-del-codice-lezione17mavendate}

In \passthrough{\lstinline!Lezione17/MavenDate!}, il prof. Pasqua
mantiene l'intera architettura pregressa e aggiunge il file
\href{file:///home/wally/Documenti/GitHub/Programmazione-II/25-26/Teoria-e-Esercitazioni/Codice-20260902/Lezione17/MavenDate/src/main/java/it/oop/ui/MainStream.java}{MainStream.java},
che esemplifica in modo magistrale tutte le funzionalità degli Stream
appena studiate.

\subsubsection{\texorpdfstring{2.1 Analisi Riga per Riga di
\href{file:///home/wally/Documenti/GitHub/Programmazione-II/25-26/Teoria-e-Esercitazioni/Codice-20260902/Lezione17/MavenDate/src/main/java/it/oop/ui/MainStream.java}{MainStream.java}}{2.1 Analisi Riga per Riga di MainStream.java}}\label{analisi-riga-per-riga-di-mainstream.java}

\begin{lstlisting}[language=Java]
package it.oop.ui;

import it.oop.core.AmericanDate;
import it.oop.core.FormattedDate;
import it.oop.core.ItalianDate;

import java.util.List;
import java.util.stream.Stream;
import it.oop.core.Date;
import static java.util.stream.Stream.generate;

public class MainStream {
    public static void main(String[] args) {
        // ESEMPIO 1: Generazione con iterate, filtro e limit
        // Sequenza generata: "a", "aa", "aaa", "aaaa", ...
        Stream<String> stream = Stream.iterate("a", s -> s + "a")
                .filter(s -> s.length() % 2 == 1) // Tiene solo stringhe di lunghezza DISPARI
                .limit(10);                       // Si arresta dopo le prime 10 stringhe dispari
        List<String> list = stream.toList();      // Operazione terminale: unwrap a List
        System.out.println(list.toString());

        // ESEMPIO 2: Generazione con generate, Supplier e test di primalità
        List<Integer> list2 = Stream.generate(() -> (int) (Math.random() * 20)) // Generatore infinito [0..19]
                .filter(i -> isPrime(i))                                        // Filtra numeri primi
                .limit(10)                                                      // Ne estrae esattamente 10
                .toList();                                                      // Operazione terminale
        System.out.println(list2.toString());

        // ESEMPIO 3: Pipeline su oggetti di dominio del progetto (Date)
        Stream.of(new ItalianDate(15,1,2025), new Date(16,1,2025), new ItalianDate(2,2,2025))
                .filter(d -> d instanceof ItalianDate) // Filtra escludendo le Date generiche
                .map(d -> new AmericanDate(d.getDay(), d.getMonth(), d.getYear())) // Converte in AmericanDate
                .forEach(System.out::println); // Method reference per stampa su stdout
    }

    // Algoritmo di primalità implementato interamente con Stream!
    private static boolean isPrime(int n) {
        return n == 0 ? false : Stream.iterate(1, i -> i + 1)
                .limit(n)              // Genera i numeri da 1 a n
                .map(i -> n % i)       // Calcola il resto della divisione (0 indica divisore esatto)
                .filter(i -> i == 0)   // Mantiene solo i resti nulli (i divisori)
                .count() <= 2;         // Se i divisori sono <= 2 (cioè 1 e se stesso), è primo!
    }
}
\end{lstlisting}

\begin{center}\rule{0.5\linewidth}{0.5pt}\end{center}

\subsection{3. Risultato di Compilazione ed Esecuzione
Reale}\label{risultato-di-compilazione-ed-esecuzione-reale}

Comando eseguito nel progetto
\href{file:///home/wally/Documenti/GitHub/Programmazione-II/25-26/Teoria-e-Esercitazioni/Codice-20260902/Lezione17/MavenDate}{Lezione17/MavenDate}:

\begin{lstlisting}[language=bash]
mvn compile exec:java -Dexec.mainClass="it.oop.ui.MainStream"
\end{lstlisting}

Output ottenuto:

\begin{lstlisting}
[a, aaa, aaaaa, aaaaaaa, aaaaaaaaa, aaaaaaaaaaa, aaaaaaaaaaaaa, aaaaaaaaaaaaaaa, aaaaaaaaaaaaaaaaa, aaaaaaaaaaaaaaaaaaa]
[7, 11, 13, 19, 3, 13, 17, 5, 19, 1]
1/15/2025
2/2/2025
\end{lstlisting}

\subsubsection{Decodifica dell'Output:}\label{decodifica-delloutput}

\begin{enumerate}
\def\labelenumi{\arabic{enumi}.}
\tightlist
\item
  \passthrough{\lstinline![a, aaa, aaaaa, ..., aaaaaaaaaaaaaaaaaaa]!}
  \(\rightarrow\) Lista delle prime 10 stringhe a lunghezza dispari (1,
  3, 5, 7, 9, 11, 13, 15, 17, 19 caratteri).
\item
  \passthrough{\lstinline![7, 11, 13, 19, 3, 13, 17, 5, 19, 1]!}
  \(\rightarrow\) Dieci numeri primi casuali estratti nell'intervallo
  \([0, 20)\) e validati dalla pipeline stream di
  \passthrough{\lstinline!isPrime(int n)!}.
\item
  \passthrough{\lstinline!1/15/2025!} e
  \passthrough{\lstinline!2/2/2025!} \(\rightarrow\) Le due istanze di
  \passthrough{\lstinline!ItalianDate!} (15/1/2025 e 2/2/2025) superano
  il \passthrough{\lstinline!filter(instanceof ItalianDate)!}, vengono
  convertite tramite \passthrough{\lstinline!map!} in
  \passthrough{\lstinline!AmericanDate!} (formato
  \passthrough{\lstinline!M/D/Y!}) e stampate a video tramite
  \passthrough{\lstinline!System.out::println!}. L'istanza generica
  \passthrough{\lstinline!new Date(16,1,2025)!} viene invece scartata
  dal filtro.
\end{enumerate}

\begin{center}\rule{0.5\linewidth}{0.5pt}\end{center}

Dimmi \textbf{``vai''} per procedere con il \textbf{Blocco 18} (Lezione
18 / Slide T19 --- \emph{Documentation and Unit Testing}, Javadoc, JUnit
3 vs JUnit 4 vs JUnit 5, asserzioni, ciclo di vita dei test
\passthrough{\lstinline!@BeforeEach!}/\passthrough{\lstinline!@AfterEach!},
e l'analisi di \passthrough{\lstinline!TestDate.java!} e
\passthrough{\lstinline!TestItalianDate.java!} in
\passthrough{\lstinline!Lezione18/MavenDate!}).


\end{document}
